\documentclass[10pt,a4paper]{article}
\usepackage[utf8]{inputenc}
\usepackage[T1]{fontenc}
\usepackage[english]{babel}
\usepackage{geometry}
\geometry{margin=1.8cm}
\usepackage{enumitem}
\usepackage{booktabs}
\usepackage{fancyhdr}
\usepackage{lastpage}
\usepackage{tcolorbox}
\usepackage{hyperref}

\pagestyle{fancy}
\fancyhf{}
\rfoot{\thepage/\pageref{LastPage}}
\lhead{\small TOEFL ITP 2026 --- Day 2}
\rhead{\small MPrometeo}
\renewcommand{\headrulewidth}{0.4pt}
\setlength{\parindent}{0pt}
\setlength{\parskip}{5pt}

\begin{document}

\begin{center}
{\LARGE \textbf{DAY 2 --- Saturday Feb 28}}\\[4pt]
{\large Listening Foundations}\\[2pt]
{\normalsize Based on Maximino Aldana:
``Yo no sab\'ia, yo no fui, era algo impredecible''}\\[2pt]
\rule{\textwidth}{0.8pt}
\end{center}

% ============ ARTICLE ============
\section*{Article Selected}

\begin{tcolorbox}[colback=white, colframe=black]
\textbf{ES:} ``Yo no sab\'ia, yo no fui, era algo impredecible''\\
\textbf{EN:} ``I Didn't Know, It Wasn't Me, It Was Unpredictable''\\
\textbf{Author:} Maximino Aldana --- La Jornada Morelos, Oct 17, 2025\\
\textbf{Link:}
\url{https://www.lajornadamorelos.mx/opinion/yo-no-sabia-yo-no-fui-era-algo-impredecible/}\\[4pt]
\textbf{Summary:} Aldana analyzes three phrases Mexican politicians
use as a manual for evading responsibility: ``I didn't know''
(claiming ignorance), ``It wasn't me'' (denying involvement), and
``It was unpredictable'' (blaming fate). He argues this pattern
extends beyond politics into Mexican society itself, normalizing
incompetence and destroying public trust. The politician who knows
nothing is as dangerous as the one who steals everything --- one
through malice, the other through incompetence.
\end{tcolorbox}

% ============ GRAMMAR ============
\section*{Grammar: Subject-Verb Agreement + Tense Basics}

\subsection*{Rule 1: Singular subject = singular verb}

Interrupting phrases (``along with,'' ``as well as,'' ``together
with'') do NOT change the subject.

\begin{itemize}[nosep]
\item The politician, along with his advisors, \textbf{denies}
      involvement. (\emph{politician} = singular)
\item Corruption, as well as impunity, \textbf{remains} systemic.
      (\emph{corruption} = singular)
\item Each of the officials \textbf{claims} innocence.
      (\emph{each} = always singular)
\end{itemize}

\subsection*{Rule 2: ``The number of'' vs.\ ``A number of''}

\begin{itemize}[nosep]
\item \textbf{The number of} scandals \textbf{has} increased.
      (singular --- the number itself)
\item \textbf{A number of} citizens \textbf{have} lost trust.
      (plural --- meaning ``many'')
\end{itemize}

\subsection*{Rule 3: Neither/nor, Either/or}

Verb agrees with the NEAREST subject:
\begin{itemize}[nosep]
\item Neither the governor \textbf{nor} his staff
      \textbf{was} able to explain.
\item Either the officials \textbf{or} the president
      \textbf{is} responsible.
\end{itemize}

\subsection*{Rule 4: Core Tenses}

\begin{center}
\begin{tabular}{lll}
\toprule
\textbf{Tense} & \textbf{Form} & \textbf{Aldana example} \\
\midrule
Simple Present & V / V+s & Corruption \textbf{destroys} trust. \\
Simple Past & V+ed / irreg & The official \textbf{denied} charges. \\
Present Perfect & has/have+pp & Impunity \textbf{has eroded} democracy. \\
Past Perfect & had+pp & He \textbf{had resigned} before the scandal. \\
Future & will+base & Citizens \textbf{will demand} accountability. \\
\bottomrule
\end{tabular}
\end{center}

% ============ 10 EXERCISES ============
\section*{Exercises: SVA + Tenses (10 questions)}

\begin{enumerate}[label=\textbf{\arabic*.}, leftmargin=*]

\item The culture of denial and impunity \rule{2cm}{0.4pt}
deeply rooted in Mexican political life.
\begin{enumerate}[label=(\Alph*)]
\item are \item is \item were \item have been
\end{enumerate}

\item Neither the president nor his cabinet members
\rule{2cm}{0.4pt} responsibility for the crisis.
\begin{enumerate}[label=(\Alph*)]
\item accepts \item accept \item accepting \item has accepted
\end{enumerate}

\item A number of journalists \rule{2cm}{0.4pt} threatened
after publishing corruption investigations.
\begin{enumerate}[label=(\Alph*)]
\item has been \item was \item have been \item is
\end{enumerate}

\item The number of disappearances in Mexico \rule{2cm}{0.4pt}
dramatically since 2006.
\begin{enumerate}[label=(\Alph*)]
\item have increased \item has increased
\item are increasing \item increase
\end{enumerate}

\item Each of the three excuses Aldana describes
\rule{2cm}{0.4pt} a different form of evasion.
\begin{enumerate}[label=(\Alph*)]
\item represent \item represents
\item are representing \item have represented
\end{enumerate}

\item By the time the scandal became public, the governor
already \rule{2cm}{0.4pt} the country.
\begin{enumerate}[label=(\Alph*)]
\item left \item has left \item had left \item leaves
\end{enumerate}

\item Public trust in institutions \rule{2cm}{0.4pt} steadily
over the past two decades.
\begin{enumerate}[label=(\Alph*)]
\item declined \item has declined
\item is declining \item decline
\end{enumerate}

\item The phrase ``it was unpredictable,'' along with
``I didn't know,'' \rule{2cm}{0.4pt} used to avoid accountability.
\begin{enumerate}[label=(\Alph*)]
\item are \item is \item were \item have
\end{enumerate}

\item If the government \rule{2cm}{0.4pt} real consequences
for corruption, the situation would improve.
\begin{enumerate}[label=(\Alph*)]
\item implements \item implemented
\item will implement \item implementing
\end{enumerate}

\item Aldana argues that politicians who claim ignorance
\rule{2cm}{0.4pt} just as dangerous as those who steal.
\begin{enumerate}[label=(\Alph*)]
\item is \item was \item are \item has been
\end{enumerate}

\end{enumerate}

\subsection*{Answer Key}

\begin{center}
\begin{tabular}{cll}
\toprule
\# & Ans & Explanation \\
\midrule
1 & B & ``The culture'' = singular $\to$ ``is'' \\
2 & B & Neither/nor: nearest = ``members'' (plural) $\to$ ``accept'' \\
3 & C & ``A number of'' = plural $\to$ ``have been'' \\
4 & B & ``The number of'' = singular $\to$ ``has increased'' \\
5 & B & ``Each'' = always singular $\to$ ``represents'' \\
6 & C & Past perfect: before another past $\to$ ``had left'' \\
7 & B & Present perfect + ``over the past'' $\to$ ``has declined'' \\
8 & B & ``The phrase, along with...'' = singular $\to$ ``is'' \\
9 & B & Second conditional: if + past $\to$ ``implemented'' \\
10 & C & ``Politicians who claim'' = plural $\to$ ``are'' \\
\bottomrule
\end{tabular}
\end{center}

% ============ READING PASSAGE ============
\section*{Reading Passage ($\sim$200 words)}

\begin{tcolorbox}[colback=white, colframe=black]
\small
In Mexico, three phrases have become an unofficial manual for
political survival: ``I didn't know,'' ``It wasn't me,'' and
``It was unpredictable.'' According to columnist Maximino Aldana,
these expressions represent more than simple excuses. They reveal
a deeply embedded cultural pattern in which accountability is
systematically avoided at every level of government.

When a corruption scandal erupts, officials typically claim
ignorance. When criminal networks are discovered operating within
state institutions, those in charge insist they had no involvement.
And when preventable tragedies occur --- flooding from poor
infrastructure, deaths from institutional negligence --- the
response is inevitably that nobody could have predicted the outcome.

Aldana argues that this cycle of denial is not limited to
politicians. It has permeated Mexican society itself, creating a
culture where responsibility is perpetually deferred. The politician
who knows nothing is as dangerous as the one who steals everything,
because both destroy public trust --- one through malice, the other
through incompetence.

Breaking this cycle, Aldana suggests, requires not only
institutional reform but a fundamental shift in how Mexican society
views accountability: not as punishment, but as the foundation of
democracy.
\end{tcolorbox}

\subsection*{Reading Questions}

\begin{enumerate}[label=\textbf{\arabic*.}]
\item What is the main idea of the passage?
\begin{enumerate}[label=(\Alph*)]
\item Mexican politicians are all corrupt
\item Three common phrases reveal systemic avoidance of accountability
\item Mexico has no democracy
\item Aldana is a politician
\end{enumerate}

\item ``Embedded'' is closest in meaning to:
\begin{enumerate}[label=(\Alph*)]
\item temporary \item deeply rooted \item visible \item foreign
\end{enumerate}

\item According to Aldana, who is more dangerous?
\begin{enumerate}[label=(\Alph*)]
\item Only corrupt politicians
\item Only ignorant politicians
\item Both are equally dangerous
\item Neither is dangerous
\end{enumerate}

\item ``Deferred'' most likely means:
\begin{enumerate}[label=(\Alph*)]
\item accepted \item postponed/passed on
\item eliminated \item celebrated
\end{enumerate}

\item What does Aldana propose?
\begin{enumerate}[label=(\Alph*)]
\item More punishment for politicians
\item Institutional reform and a cultural shift toward accountability
\item Replacing all officials
\item Ignoring the problem
\end{enumerate}
\end{enumerate}

\textbf{Answers:} 1-B, 2-B, 3-C, 4-B, 5-B

% ============ SHADOWING ============
\section*{Shadowing Exercise (15 min --- Gemini Live)}

\begin{tcolorbox}[colback=white, colframe=black,
    title=Tell Gemini Live:]
``Read the following paragraph aloud at normal American English
speed. Then I will repeat it sentence by sentence. Correct my
pronunciation after each sentence.''
\end{tcolorbox}

\begin{quote}\small
``In Mexico, three phrases have become an unofficial manual for
political survival. When a corruption scandal erupts, officials
typically claim ignorance. When criminal networks are discovered
operating within state institutions, those in charge insist they
had no involvement. This cycle of denial has permeated Mexican
society itself, creating a culture where responsibility is
perpetually deferred. Breaking this cycle requires a fundamental
shift in how society views accountability.''
\end{quote}

\textbf{Pronunciation targets:}
\begin{itemize}[nosep]
\item \textbf{corruption} --- stress on 2nd syllable: cor-RUP-tion
\item \textbf{accountability} --- stress on 3rd: ac-count-A-bil-i-ty
\item \textbf{permeated} --- stress on 1st: PER-me-a-ted
\item \textbf{perpetually} --- stress on 2nd: per-PETCH-u-al-ly
\item \textbf{inevitably} --- stress on 2nd: in-EV-i-ta-bly
\end{itemize}

\end{document}